\chapter{The isoperimetric equality case}\label{ch:rigidity}
\module{NoCompromise.Isoperimetric.Rigidity}

\begin{theorem}[Isoperimetric rigidity]\label{thm:isoperimetric-rigidity}
\uses{thm:sharp-isoperimetric,prop:no-singular-points}
Let $E\subset\R^3$ have $0<|E|<\infty$, finite perimeter, and
$\Per(E)=c_I|E|^{2/3}$.  Then $E$ is, up to a null set, a translate of a ball
of volume $|E|$.  Conversely balls give equality.
\end{theorem}

\begin{proof}
\uses{lem:penalization,lem:quasiminimal,prop:no-singular-points,
      cor:iso-smooth-components,lem:concavity-23,lem:abp-neumann,
      lem:abp-contact,thm:boundary-neumann,prop:iso-smooth,
      lem:ball-perimeter}
\emph{Step 1: $E$ is $\omega$-minimal.}  Equality makes $E$ a perimeter
minimiser among sets of its volume.  Apply the proof of
Lemma~\ref{lem:penalization} with the Coulomb term deleted --- i.e.\ with
$\Ec$ replaced by $\Per$ --- which gives a finite volume-penalty constant, and
then the proof of Lemma~\ref{lem:quasiminimal} with the Coulomb Lipschitz term
deleted.  Thus the generic $\omega$-minimal regularity theory applies with no
liquid-drop term built into it.

\emph{Step 2: a bounded smooth representative.}
Chapters~\ref{ch:density}--\ref{ch:no-singular} apply verbatim and give a
bounded representative $\Omega$ with compact $C^{1,1/2}$ boundary; the
bootstrap of \S\ref{sec:bootstrap} with $v_\Omega$ replaced by $0$ makes it
$C^{3,\beta}$, which is all that Steps 3--5 need: the ABP chain of
Proposition~\ref{prop:iso-smooth} and Corollary~\ref{cor:iso-smooth-components}
uses only a $C^3$ boundary (see the formalisation note after
Proposition~\ref{prop:iso-smooth}).  \emph{These chapters use only the inequality
\eqref{eq:sharp-isoperimetric}, never this theorem.}

\emph{Step 3: one component.}  Apply
Corollary~\ref{cor:iso-smooth-components}; by
Lemma~\ref{lem:concavity-23} equality leaves exactly one component.

\emph{Step 4: equality in the ABP chain.}  In the chain of
Proposition~\ref{prop:iso-smooth} every integrand inequality is nonnegative,
so equality of the endpoints implies
\begin{itemize}[nosep]
\item $|\nabla z(\Gamma)\setminus B_1|=0$;
\item $|G\setminus\Gamma|=0$, because $\Delta z/3>0$;
\item $D^2z=(\Delta z/3)I$ a.e.\ on $G$, by equality in the
  arithmetic--geometric mean inequality for a positive semidefinite matrix.
\end{itemize}

\emph{Step 5: conclude.}  Smoothness extends the Hessian identity everywhere,
so $\nabla z(x)=\alpha x+b$ with $\alpha=\Per(E)/(3|E|)>0$.  The affine image
$\nabla z(G)$ is open, contains $B_1$, and has the same volume as $B_1$; an
open proper superset would have strictly larger volume, so
$\nabla z(G)=B_1$ and $G=\alpha^{-1}(B_1-b)$ is a ball.  The converse is
Lemma~\ref{lem:ball-perimeter} and a direct computation.
\end{proof}

